\HLASourcePage{26}{25}
\thispagestyle{hla-body}
\begin{hlafixedpage}
  \HLADefinitionAt{69.79}{%
    \HLADefinitionAnchor{delete}\HLADefinitionTerm{delete}: To invoke the \emph{Delete Object Instance} service to eliminate a particular object instance. \emph{See also}:\linebreak
    \textbf{remove}.%
  }

  \HLADefinitionNoteAt{90}{105.98}{NOTE---See 6.16.}

  \HLADefinitionAt{127.81}{%
    \HLADefinitionAnchor{designator}\HLADefinitionTerm{designator}: A generic view of arguments referenced in service descriptions. This view improves clarity in\linebreak
    situations when arguments (mostly identifiers) to services have different views (implementations) due to a\linebreak
    particular programming language or application programmer's interface (API).%
  }

  \HLADefinitionAt{175.81}{%
    \HLADefinitionAnchor{dimension}\HLADefinitionTerm{dimension}: A named interval. The interval is defined by an ordered pair of values, the first being the\linebreak
    dimension lower bound and the second being the dimension upper bound. A runtime infrastructure (RTI)\linebreak
    provides a predefined interval, whose lower and upper bounds are fixed as [0, the dimension's upper bound)\linebreak
    as specified in the federation object model (FOM) Document Data (FDD). This interval provides a single\linebreak
    basis for communication of all dimension-related data with an RTI. All normalized intervals communicated\linebreak
    to the RTI will be subsets of this interval.%
  }

  \HLADefinitionNoteAt{90}{260.00}{NOTE---See 9.1.}

  \HLADefinitionAt{281.77}{%
    \HLADefinitionAnchor{directed-interaction}\HLADefinitionTerm{directed interaction}: An explicit action taken by a federate directed to a specific object instance that may\linebreak
    have some effect or impact on another federate within a federation execution that owns attributes of that\linebreak
    object instance.%
  }

  \HLADefinitionAt{329.76}{%
    \HLADefinitionAnchor{directed-interaction-class}\HLADefinitionTerm{directed interaction class}: An association between an interaction class and an object class where instances\linebreak
    of the interaction class may be directed to be received by interested federate that owns attribute instances of\linebreak
    the object class.%
  }

  \HLADefinitionAt{375.57}{%
    \HLADefinitionAnchor{discover}\HLADefinitionTerm{discover}: To receive an invocation of the \emph{Discover Object Instance} \textdagger{}\HLASupMark{13} service for a particular object\linebreak
    instance.%
  }

  \HLADefinitionNoteAt{90}{413.96}{NOTE---See 6.9.}

  \HLADefinitionAt{435.79}{%
    \HLADefinitionAnchor{discovered-class}\HLADefinitionTerm{discovered class}: The class that was an object instance's candidate discovery class at a joined federate when\linebreak
    that object instance was discovered by that joined federate. \emph{See also}: \textbf{candidate discovery class}.%
  }

  \HLADefinitionNoteAt{90}{471.98}{NOTE---See 5.1.2.}

  \HLADefinitionAt{493.81}{%
    \HLADefinitionAnchor{dynamic-link-compatibility}\HLADefinitionTerm{dynamic link compatibility (DLC)}: The ability to interchange the implementation of an API with another\linebreak
    without having to recompile or relink the application software that uses the API.%
  }

  \HLADefinitionAt{529.81}{%
    \HLADefinitionAnchor{epsilon-time-interval}\HLADefinitionTerm{epsilon time interval}: The distance between any two consecutive logical times T1 and T2 (i.e. the set of\linebreak
    logical times between T1 and T2 is empty). For some representations of local time (e.g. floating point), the\linebreak
    distance may not be the same for all such instances of T1 and T2.%
  }

  \HLADefinitionAt{577.80}{%
    \HLADefinitionAnchor{exception}\HLADefinitionTerm{exception}: Notification of any irregularity that may occur during a service invocation.%
  }

  \HLADefinitionAt{601.80}{%
    \HLADefinitionAnchor{fault}\HLADefinitionTerm{fault}: The occurrence of an unexpected event, at any point during a federation execution, that prevents one\linebreak
    or more joined federates in the federation from executing in an HLA-compliant manner.%
  }

  \HLADefinitionAt{637.80}{%
    \HLADefinitionAnchor{federate}\HLADefinitionTerm{federate}: An application that may be or is currently coupled with other software applications under a\linebreak
    federation object model (FOM) Document Data (FDD) and a runtime infrastructure (RTI). \emph{See also}:\linebreak
    \textbf{federate application}; \textbf{joined federate}.%
  }

  \draw[line width=0.48pt] (90.78,701.90)--(127.26,701.90);
  \HLAFixedBlock{90}{704.10}{6}{\rmfamily\raggedright}{6.36}{7}{13}
  \HLAFixedBlock{96.36}{707.25}{425.64}{\rmfamily\raggedright}{7.98}{9.4}{%
    All RTI-initiated services are denoted with a \textdagger{} (printer's dagger) after the service name.%
  }
\end{hlafixedpage}
\clearpage
